% Синтетический кейс: теоремы через thmtools (\declaretheorem), номера из .aux.
\documentclass[12pt]{article}
\usepackage[T2A]{fontenc}
\usepackage[utf8]{inputenc}
\usepackage[russian]{babel}
\usepackage{amsmath,amsthm,thmtools}

\declaretheorem[name=Теорема, numberwithin=section]{theorem}
\declaretheorem[name=Лемма, sibling=theorem]{lemma}
\declaretheorem[name=Определение, style=definition, sibling=theorem]{definition}
\declaretheorem[name=Замечание, style=remark, numbered=no]{remark}

\begin{document}

\section{Теоремы thmtools}

\begin{definition}\label{def:t}
Определение через thmtools.
\end{definition}

\begin{theorem}[Главная]\label{thm:t}
Теорема с именем.
\end{theorem}

\begin{proof}
По определению~\ref{def:t}.
\end{proof}

\begin{lemma}
Лемма без метки.
\end{lemma}

\begin{remark}
Замечание без номера.
\end{remark}

\end{document}
